% =============================================================================
% §5.7 (or insert into §4.4) — Exemplar-correction re-experiment.
% Self-contained: no new commands, no new packages required beyond those
% already loaded in main_journal.tex (booktabs, amsmath).
% =============================================================================

\subsection{Isolating exemplar quality from prompting strategy}
\label{sec:exemplar-correction}

A reviewer correctly noted that the headline structured-vs.-few-shot gap
(10/10 vs.\ 1/10) is partially confounded by the few-shot exemplar itself:
the version of \texttt{wsi\_mlp\_train\_fl\_client.py} used to generate the
cached \texttt{\_fl\_few\_shot.py} outputs contained an I3 contract
violation in \texttt{train\_step} (\texttt{loss.backward();
optimizer.step(); return loss.detach()}), and a static AST audit confirms
that 5/10 of those cached outputs copied the violation verbatim. To
disentangle exemplar quality from prompting strategy, we re-ran the
few-shot condition with a contract-correct exemplar (added as
\texttt{Skill.FEW\_SHOT\_CORRECTED} in \texttt{converter/llm\_converter.py})
in a dry-run analysis: each cached few-shot output had its
contract-violating lines mechanically removed (drop \texttt{.backward()},
\texttt{optimizer.step()}, \texttt{optimizer.zero\_grad()},
\texttt{scaler.\{scale,step,update,unscale\_\}()}; rewrite
\texttt{return X.detach()/X.item()} to \texttt{return X}) and was
re-evaluated through the unchanged preflight + 1-round FL pipeline; this
yields a sound \emph{upper bound} on the corrected-exemplar success rate
because the patch strictly removes violations while preserving every other
LLM choice. The corrected-exemplar few-shot reached \textbf{4/10}
end-to-end success---a 30-point absolute gain over the original 1/10---
with all three newly-passing scripts (\texttt{dcgan\_main},
\texttt{mnist\_main}, \texttt{backbone\_image\_classifier}) being precisely
those whose only failure mode was the propagated I3 violation; the
remaining six failures (missing \texttt{keras}/\texttt{nibabel} modules,
missing dataset directories, a MONAI empty-dataset path) are environmental
and would not be cured by any exemplar change. Two-sided Fisher's exact
tests give Original-vs-Corrected $p = 0.30$ (not significant at
$n=10$, although the 30-point point estimate is large and
mechanism-consistent), Corrected-vs-Structured $p = 0.011$, and
Original-vs-Structured $p = 1.2 \times 10^{-4}$. We therefore conclude
that exemplar quality and prompting strategy are co-equal factors:
roughly one-third of the original 10$\times$ structured-vs.-few-shot gap
is attributable to a single bad demonstration line in the exemplar, but
the residual 4/10-vs.-10/10 gap---which corresponds to scripts that need
behavioural specifications no positive example can carry (synthetic-data
fallbacks, optional-module guards, contract \emph{prohibitions})---remains
significant and supports the paper's central claim that structured
specifications outperform demonstration-based prompting on tasks dominated
by negative behavioural constraints.
